\documentclass[10pt,a4paper]{amsart}
\usepackage[margin=19mm]{geometry}
\usepackage{amsmath,amssymb,mathtools}
\usepackage[scr=rsfs,cal=euler]{mathalfa}
\usepackage{xfrac}
\usepackage[unicode,hidelinks,bookmarks=false]{hyperref}
\newcommand{\ie}{i.e.\ }
\newcommand{\cf}{cf.\ }
\newcommand{\resp}{respectively\ }
\newcommand{\Subsection}{\subsection}
\providecommand{\nobreakdash}{\nobreak\hbox{-}}
\setlength{\emergencystretch}{2em}
\multlinegap0pt
\usepackage{amssymb}
\usepackage{mathtools}
\usepackage[shortlabels]{enumitem}
\newtheorem{lemma}{Lemma}
\newtheorem{proposition}{Proposition}
\newcommand{\B}{\mathbb B}
\newcommand{\Rand}{\operatorname{Rand}}
\newcommand{\restr}{\mathbin{\upharpoonright}}
\title{Equivalence of Random Generics Is Not Essentially Free}
\author{Jason Zesheng Chen$^\ast$}
\date{Source: arXiv:2608.21706v1 (CC BY 4.0)}
\begin{document}
\maketitle

\noindent\emph{Source and licence.} Equivalence of Random Generics Is Not Essentially Free. Version arXiv:2608.21706v1 (CC BY 4.0), distributed by arXiv under the Creative Commons Attribution 4.0 International licence, http://creativecommons.org/licenses/by/4.0/, as recorded in the arXiv licence metadata for this exact version and shown on the abstract page https://arxiv.org/abs/2608.21706v1. Full licence notice: \texttt{licenses/arxiv-2608.21706-CC-BY-4.0.txt}.\par\smallskip

\noindent\textit{Editorial scope.} Counted unit: the article's proposition prop:second-coordinate (source0.tex lines 302--314). The article's complete original proof of it is delivered with it (source0.tex lines 316--347, one \texttt{proof} environment of 32 lines). No mathematics is written, repaired or re-derived: the statement and the proof are the article's own lines, in the article's own notation, and no retained line is substituted or re-flowed.  Retained uncounted as the source-local context the proof needs: lines 144--149.  Nothing was omitted from any retained span, so the package carries no omission record.  No retained line contains a citation, so the wrapper carries no bibliography and no bibliography entry.  No retained line contains a figure environment, an image-inclusion command or drawing code, so no image is delivered and the package declares no asset dependency.  Editorial formatting only, in addition: an AMS \texttt{amsart} preamble that restates the article's own declarations and macros used by the retained lines, namely the environments \texttt{lemma}, \texttt{proposition} on separate counters, together with the standard packages those lines need.  The retained environments print in this document as Lemma 1, Proposition 1 in order of appearance, whereas the article prints the same items with the numbering of its own full document.  The complete native source of the article as distributed in the arXiv e-print for this exact version is bundled unchanged as \texttt{sources/208271/source0.tex}.
\begin{lemma}\label{lem:two-step}
Suppose that $x$ is random over $N$ and $u$ is random over $N[x]$.  Then $(x,u)$ is product-random over $N$.  Consequently $I(x,u)$ is random over $N$ and
\[
 N[I(x,u)]=N[x,u].
\]
\end{lemma}

\begin{proposition}\label{prop:second-coordinate}
Let $N$ be a countable transitive model of a sufficiently large finite fragment of ZFC, let $x$ be random over $N$, and let $\Delta\in N[x]$ be a group which $N[x]$ regards as countably infinite.  Suppose that $C\subseteq2^\omega$ is Borel and $\mu(C_x)=1$.  Then there are an invariant conull Borel set $R\subseteq2^\Delta$ and an injective Borel map
\[
 c\colon R\longrightarrow C\cap\Rand(N)
\]
such that
\begin{enumerate}
\item $R$ is contained in the free part of the Bernoulli shift $\Delta\curvearrowright2^\Delta$;
\item $N[c(z)]=N[x,z]$ for every $z\in R$;
\item $c(z)\,E_N^{\B}\,c(\delta\cdot z)$ for every $z\in R$ and $\delta\in\Delta$.
\end{enumerate}
Consequently, there is a Borel homomorphism from $E(\Delta,2)\restr R$ to $E_N^{\B}\restr(C\cap\Rand(N))$.
\end{proposition}

\begin{proof}
Choose in $N[x]$ a bijection $e\colon\omega\to\Delta$ and let
\[
 e^*\colon2^\Delta\longrightarrow2^\omega,
 \qquad e^*(z)(n)=z(e(n)),
\]
be the induced measure-preserving homeomorphism.  Put
\[
 D=\{z\in2^\Delta:I(x,e^*(z))\in C\}.
\]
Since $\mu(C_x)=1$, the set $D$ is conull.  Put
\[
 D_0=\bigcap_{\delta\in\Delta}\{z\in2^\Delta:\delta\cdot z\in D\}.
\]
Since $\Delta$ is countable, $D_0$ is Borel, invariant, and conull.

Let $R$ be the intersection of $D_0$ with the set of points of $2^\Delta$ which are random over $N[x]$ for Bernoulli product measure.  Externally $N[x]$ is countable, so $R$ is Borel and conull.  It is invariant because every shift and its inverse are measure-preserving Borel maps coded in $N[x]$.  The free part of the Bernoulli action is an $N[x]$-coded conull Borel set, and therefore every point random over $N[x]$ belongs to it.

For $z\in R$, define
\[
 c(z)=I(x,e^*(z)).
\]
The map $c$ is Borel and injective, and $c(z)\in C$ because $z\in D$.  Since $e^*$ is a measure-preserving Borel isomorphism coded in $N[x]$, the real $e^*(z)$ is random over $N[x]$.  By Lemma~\ref{lem:two-step}, $c(z)$ is random over $N$.  Since $I^{-1}$ recovers $x$ and $e^*(z)$, while $e\in N[x]$, one has
\[
 N[c(z)]=N[x,e^*(z)]=N[x,z].
\]
For $\delta\in\Delta$, the shift $z\mapsto\delta\cdot z$ and its inverse belong to $N[x]$.  Hence
\[
 N[x,z]=N[x,\delta\cdot z],
\]
and therefore $N[c(z)]=N[c(\delta\cdot z)]$.
\end{proof}

\end{document}
